\chapter*{กิตติกรรมประกาศ}
\addcontentsline{toc}{chapter}{กิตติกรรมประกาศ}

เอกสารคำสอนและคู่มือปฏิบัติการรายวิชา 4011111 ฟิสิกส์สำหรับจุลชีววิทยา เล่มนี้ สำเร็จลุล่วงได้อย่างสมบูรณ์ด้วยความอนุเคราะห์และการสนับสนุนทางวิชาการจากคณะวิทยาศาสตร์และเทคโนโลยี มหาวิทยาลัยราชภัฏรำไพพรรณี ที่ได้ส่งเสริมการพัฒนาสื่อการสอนและเอกสารคำสอนที่สอดคล้องกับเกณฑ์มาตรฐานหลักสูตรและการจัดการศึกษาที่มุ่งเน้นผลลัพธ์การเรียนรู้

ผู้เขียนขอกราบขอบพระคุณคณาจารย์ประจำหลักสูตรวิทยาศาสตรบัณฑิต สาขาวิชาจุลชีววิทยา และคณาจารย์สาขาวิชาฟิสิกส์ทุกท่าน ที่ได้ร่วมแลกเปลี่ยนองค์ความรู้และให้ข้อเสนอแนะในการคัดเลือกหัวข้อเนื้อหาและออกแบบชุดปฏิบัติการทดลองให้เชื่อมโยงโดยตรงกับเครื่องมือวิทยาศาสตร์ในห้องปฏิบัติการจุลชีววิทยา อาทิ เครื่องหมุนเหวี่ยงตกตะกอน หม้อนึ่งความดันไอ ตู้อบเลี้ยงเชื้อ ชุดอิเล็กโทรโฟรีซิส กล้องจุลทรรศน์กำลังขยายสูง และเครื่องสเปกโตรโฟโตมิเตอร์

ขอขอบคุณนักศึกษาสาขาวิชาจุลชีววิทยาทุกชั้นปี ที่ได้สะท้อนความคิดเห็นเกี่ยวกับการเรียนรู้วิชาฟิสิกส์พื้นฐาน ซึ่งเป็นแรงบันดาลใจสำคัญให้ผู้เขียนพัฒนาวิธีอธิบายด้วยภาพจำสามเหลี่ยมสูตรคำนวณและขั้นตอนการแทนค่าที่เข้าใจง่าย ช่วยลดอุปสรรคด้านคณิตศาสตร์และเพิ่มความสนุกในการทดลองจริง

คุณค่าและประโยชน์ทางวิชาการอันพึงมีจากเอกสารคำสอนเล่มนี้ ผู้เขียนขอน้อมบูชาพระคุณบิดา มารดา และบูรพาจารย์ทุกท่านที่ได้ประสิทธิ์ประสาทวิชาความรู้และเป็นแบบอย่างของความเป็นครูผู้มุ่งมั่นพัฒนาผู้เรียนเสมอมา

\vspace{1.2cm}
\begin{flushright}
    \textbf{ผู้ช่วยศาสตราจารย์ ดร.ชีวะ ทัศนา}\\[0.2cm]
    สาขาวิชาฟิสิกส์ คณะวิทยาศาสตร์และเทคโนโลยี\\
    มหาวิทยาลัยราชภัฏรำไพพรรณี
\end{flushright}
